% !TEX program = lualatex
% Auto-generated from YAML (v3) — 13: デジタル表現1：2進数 / 文字コード / 音

\documentclass[aspectratio=43,professionalfonts,handout,10pt]{beamer}
\usepackage{beamer_template}

\newcommand{\LectureNum}{13}
\newcommand{\LectureTitle}{デジタル表現1：2進数 / 文字コード / 音}
\newcommand{\CourseName}{情報リテラシー}
\newcommand{\TextPages}{15-17}

\begin{document}
% --- Slide 1: title ---

\begin{frame}{\CourseName\quad 第\LectureNum 回「\LectureTitle 」}
  {\small \TermName\quad \TeacherName\quad ／\quad 教科書 \TextPages}

  \vfill

  \begin{block}{今日の流れ}
    \begin{enumerate}
      \item 2進法の計算・補数（Theory 15）
      \item 文字コードの仕組み（Theory 16）
      \item 音のデジタル化（Theory 17）
    \end{enumerate}
  \end{block}
  \vfill

  \begin{block}{今日の目標}
    \begin{itemize}
      \item 補数の仕組み，文字コード，音声のデジタル化について説明できる。
    \end{itemize}
  \end{block}
  \vfill

  \begin{exampleblock}{キーワード}
    \small \textbf{文字コード} ／ \textbf{ASCII} ／ \textbf{JIS8ビットコード} ／ \textbf{シフトJISコード} ／ \textbf{ユニコード（Unicode）} ／ \textbf{文字化け} ／ \textbf{振幅} ／ \textbf{周波数（Hz）} ／ \textbf{標本化（サンプリング）} ／ \textbf{量子化} ／ \textbf{符号化}
  \end{exampleblock}
\end{frame}

% --- Slide 2: 補数 ---

\begin{frame}{補数}
  \begin{block}{ポイント}
    \begin{itemize}
      \item 補数とは，ある数に足すと基数（2進法なら2）のべき乗になる数のこと
      \item コンピュータでは負の数を表すために補数を利用する
      \item 2進nビットの補数：元の数のビットを反転して1を加える
      \item 例：4ビットの1011の補数 → 0100 + 1 = 0101（10−11=−1の計算に利用）
    \end{itemize}
  \end{block}

  \vfill

  \begin{exampleblock}{補足}
    \begin{itemize}
      \item 補数を使うと，引き算を足し算として計算できる（コンピュータの演算の基本）
    \end{itemize}
  \end{exampleblock}

  \begin{slidenote}
    黒板で手順を示す：1011の補数 → ①ビット反転：0100 ②1を加える：0101。確認：$1011+0101=10000$（4ビットでは繰り上がりを無視して0000）→差が0。これが「引き算を足し算で実現」できる仕組み。符号付き4ビットでは最上位ビットが符号ビット（0=正，1=負）。
  \end{slidenote}
\end{frame}

% --- Slide 3: 文字コードの仕組み ---

\begin{frame}{文字コードの仕組み}
  \begin{block}{ポイント}
    \begin{itemize}
      \item 文字コード：文字に番号（コード）を割り当てて2進数で表現する仕組み
      \item ASCII：英数字・記号を7ビット（128文字）で表す基本的な文字コード
      \item JIS8ビットコード：ASCIIを8ビットに拡張し，カタカナを追加
      \item シフトJISコード：JISコードを改良し，漢字・ひらがなを2バイトで表現
      \item ユニコード（Unicode）：世界中の文字を統一的に扱える国際標準の文字コード
    \end{itemize}
  \end{block}

  \vfill

  \begin{exampleblock}{補足}
    \begin{itemize}
      \item 文字化け：送り手と受け手で異なる文字コードを使うと発生する
    \end{itemize}
  \end{exampleblock}

  \begin{slidenote}
    ASCIIは7ビット（128文字）で英語圏向け。日本語対応でシフトJIS（漢字・ひらがなを2バイト）が生まれ，現在はWebの標準がUTF-8（Unicode）。Unicodeはアラビア語・絵文字まで含む14万字以上に対応。文字化けの体験談：メールの件名が「?????」になった経験を聞いてみる。
  \end{slidenote}
\end{frame}

% --- Slide 4: 音のデジタル化 ---

\begin{frame}{音のデジタル化}
  \begin{block}{ポイント}
    \begin{itemize}
      \item 音はアナログ（空気の連続的な振動）→ A/D変換でデジタル化
      \item 音を特徴づける要素：振幅（音の大きさ）・周波数（音の高低）・波形（音色）
      \item ①標本化（サンプリング）：一定時間間隔で波の高さを取り出す
      \item ②量子化：取り出した値を一定の段階値に丸める
      \item ③符号化：量子化した数値を2進法で表す
    \end{itemize}
  \end{block}

  \begin{slidenote}
    CDの規格を例に：標本化44,100Hz（1秒間に44,100回），量子化16ビット（65,536段階）。人間の可聴域（20〜20,000Hz）をカバーするにはナイキスト定理より2倍以上の周波数が必要→44,100Hz。①標本化→②量子化→③符号化の流れを黒板で図示する。
  \end{slidenote}
\end{frame}

% --- Slide 5: 音のデータ量計算 ---

\begin{frame}{音のデータ量計算}
  \begin{block}{ポイント}
    \begin{itemize}
      \item 1秒間のデータ量 = 標本化周波数（Hz）× 量子化ビット数（bit）× チャンネル数
      \item 例：標本化周波数24,000Hz，量子化16bit，ステレオ（2ch）の場合
      \item 24,000 × 16 × 2 = 768,000 bit/秒（= 768 kbit/秒）
      \item 品質↑（標本化周波数・量子化ビット数を増やす）→ データ量↑（トレードオフ）
    \end{itemize}
  \end{block}

  \begin{slidenote}
    黒板で計算：$24{,}000\times16\times2=768{,}000$ bit/秒。1分間では$768{,}000\times60=46{,}080{,}000$ bit $\approx5.49$ MB。CDの場合（$44{,}100\times16\times2$）≈1.4Mbit/秒，1分≈10MB。品質とデータ量のトレードオフは演習の前振りとして強調する。
  \end{slidenote}
\end{frame}

% --- Slide 6: 演習 ---

\begin{frame}{演習}
  {\small 各自で取り組んでください。}

  \vfill

  \begin{block}{演習1}
    4ビットの2進数 0110 の補数を求めよ
  \end{block}
  \vfill

  \begin{block}{演習2}
    ASCIIとユニコードの違いを説明せよ
  \end{block}
  \vfill

  \begin{block}{演習3}
    「標本化」「量子化」「符号化」を，音のデジタル化を例にして説明せよ
  \end{block}

  \begin{slidenote}
    解答例：演習1：0110を反転→1001，$+1$→1010。演習2：ASCIIは英数字128文字のみ（7ビット），Unicodeは世界中の文字を統一的に表現（14万字以上，UTF-8等で可変長エンコード）。演習3：音の波形を一定間隔で取り出す（標本化）→段階値に丸める（量子化）→2進数に変換（符号化）。
  \end{slidenote}
\end{frame}

% --- チェックリスト ---

\begin{frame}{まとめ・チェックリスト}
  \begin{block}{自己チェック（できたら $\checkmark$ をつけよう）}
    \begin{itemize}
      \item[$\square$] 2進法の加算と補数（負の数の表現）を習得した
      \item[$\square$] 文字コード（ASCII・Unicode）の概念と文字化けの原因を理解した
      \item[$\square$] 音のデジタル化（標本化・量子化・符号化）とデータ量計算を学んだ
    \end{itemize}
  \end{block}

  \vfill

  {\small 質問があれば授業後またはTeamsで受け付けます。}
\end{frame}

\end{document}
